\documentclass[11pt]{article}

\usepackage[margin=1in]{geometry}
\usepackage{amsmath,amssymb,amsthm,mathtools}
\usepackage{enumitem}
\usepackage{hyperref}

\newtheorem{theorem}{Theorem}
\newtheorem{lemma}{Lemma}
\newtheorem{corollary}{Corollary}
\newtheorem{remark}{Remark}

\DeclareMathOperator{\OPT}{OPT}
\newcommand{\E}{\mathbb{E}}
\newcommand{\Prb}{\mathbb{P}}
\newcommand{\1}{\mathbf{1}}
\newcommand{\dd}{\mathrm{d}}
\newcommand{\cost}{\operatorname{cost}}

\title{A Degree-Constrained LP Witness for the Backbone Circulation}
\author{}
\date{}

\begin{document}
\maketitle

\section{Setting}

Let $(\{r\}\cup V,d)$ be a directed metric, where $r$ is the depot.  Put
$N:=\{r\}\cup V$ and $n:=\max\{|N|,3\}$, so that all logarithms below are
at least a positive absolute constant.  Each
$v\in V$ is active independently with probability $p_v$.  For an active set
$A\subseteq V$, let $\OPT(A)$ denote the length of an optimal directed tour on
$\{r\}\cup A$.  Define
\[
  \OPT_{\mathrm{post}} := \E_A[\OPT(A)].
\]

For a set $S\subseteq V$, write $U:=V\setminus S$.  We consider the following
backbone circulation LP with an additional degree bound $D$ on the backbone:
\[
\begin{aligned}
\min \quad & \sum_{u,v\in \{r\}\cup V} d(u,v)x_{uv} \\
\text{s.t.}\quad
& x(\delta^+(w))=x(\delta^-(w)) && \forall w\in \{r\}\cup V, \\
& x(\delta^+(v))=x(\delta^-(v))\ge p_v && \forall v\in U, \\
& x(\delta^+(s))=x(\delta^-(s))\le D && \forall s\in S\cup\{r\}, \\
& x_{uv}=0 && \forall u,v\in U, \\
& x_{uv}\ge 0 && \forall u,v.
\end{aligned}
\tag{$\mathrm{LP}_D(S)$}
\]

The goal of this note is to prove that the additional constraints
\[
  x(\delta^-(s))\le O(\log n), \qquad s\in S\cup\{r\},
\]
do not destroy the comparison with $\OPT_{\mathrm{post}}$.  More precisely,
we prove that there exists a feasible solution satisfying the degree constraints
whose cost is at most $O(\log n)\OPT_{\mathrm{post}}$.

\section{The two-sample witness}

Let $B\subseteq V$ be an ordinary active sample, i.e. each $v\in V$ is included
in $B$ independently with probability $p_v$.  We always regard $r$ as part of
the backbone, so the backbone is $B\cup\{r\}$.

Let $A\subseteq V$ be an independent active sample, and define
\[
  C := A\cup B.
\]
For every possible set $C$, fix once and for all an optimal directed tour
$T_C$ on $\{r\}\cup C$.  The tour $T_C$ depends only on $C$, not on the
particular pair $(A,B)$ that generated $C$.

Fix a pair $(A,B)$.  Cut $T_C$ at the vertices of $B\cup\{r\}$.  This gives
cyclically ordered directed segments
\[
  P_j=(b_j,u_{j,1},u_{j,2},\ldots,u_{j,k_j},b_{j+1}),
\]
where $b_j,b_{j+1}\in B\cup\{r\}$ and all internal vertices
$u_{j,t}$ lie in $A\setminus B$.  Let
\[
  M(A,B):=\max_j k_j.
\]

We now define a circulation $y^{A,B}$.  For every segment
$P_j=(b_j,u_{j,1},\ldots,u_{j,k_j},b_{j+1})$:
\begin{itemize}[leftmargin=2em]
  \item for each internal vertex $u_{j,t}$, send one unit of flow on
  $b_j\to u_{j,t}$ and one unit of flow on $u_{j,t}\to b_{j+1}$;
  \item send $M(A,B)-k_j$ units of flow on the backbone arc
  $b_j\to b_{j+1}$.
\end{itemize}

\begin{lemma}[Circulation and degrees for a fixed pair]
For every fixed pair $(A,B)$, $y^{A,B}$ is a circulation.  Every vertex in
$A\setminus B$ has one unit of incoming flow and one unit of outgoing flow.
Every vertex in $B\cup\{r\}$ has $M(A,B)$ units of incoming flow and
$M(A,B)$ units of outgoing flow.
\end{lemma}

\begin{proof}
Fix a segment $P_j$.  The total flow leaving its first backbone endpoint
$b_j$ through this segment is
\[
  k_j+(M(A,B)-k_j)=M(A,B).
\]
Similarly, the total flow entering the next backbone endpoint $b_{j+1}$ from
this segment is also $M(A,B)$.  Since the segments form a cyclic decomposition
of $T_C$, every backbone vertex has exactly $M(A,B)$ incoming flow and
$M(A,B)$ outgoing flow.  Each internal vertex $u_{j,t}\in A\setminus B$ receives
one unit from $b_j$ and sends one unit to $b_{j+1}$.  Hence all vertices are
balanced.
\end{proof}

Now average over the independent sample $A$, while keeping $B$ fixed:
\[
  x^B := \E_A[y^{A,B}\mid B].
\]

\begin{lemma}[Feasibility before the degree bound]
For every fixed $B$, the vector $x^B$ is feasible for the backbone LP on
$B\cup\{r\}$ without the upper degree constraint.  Moreover, for every
$v\notin B$,
\[
  x^B(\delta^+(v))=x^B(\delta^-(v))=p_v,
\]
and for every $b\in B\cup\{r\}$,
\[
  x^B(\delta^+(b))=x^B(\delta^-(b))=
  Z(B),
  \qquad
  Z(B):=\E_A[M(A,B)\mid B].
\]
\end{lemma}

\begin{proof}
The average of circulations is a circulation.  No arc between two vertices of
$V\setminus B$ is ever used: every used arc is either from $B\cup\{r\}$ to an
internal vertex, from an internal vertex to $B\cup\{r\}$, or between two
vertices of $B\cup\{r\}$.  Hence the support constraint is satisfied.

Fix $v\notin B$.  For a given $A$, the vertex $v$ is internal in exactly one
segment if $v\in A$, and is absent otherwise.  Thus $y^{A,B}$ gives $v$ one
unit of incoming and outgoing flow exactly on the event $v\in A$.  Averaging
over $A$ gives degree $p_v$ at $v$.

The formula for backbone degrees follows directly from the previous lemma:
for every fixed $A$, each backbone vertex has incoming and outgoing degree
$M(A,B)$; averaging gives $Z(B)$.
\end{proof}

\section{The logarithmic degree bound}

We next show that $Z(B)=O(\log n)$ with high probability over $B$.

\begin{lemma}[Longest run bound]
There is an absolute constant $c_0$ such that, for every $t\ge 1$,
\[
  \Prb[M(A,B)\ge t]\le n2^{-t}.
\]
Consequently,
\[
  \E[M(A,B)^q]^{1/q}\le c_0(\log n+q)
\]
for every integer $q\ge 1$.
\end{lemma}

\begin{proof}
Condition on $C=A\cup B$.  Then the tour $T_C$ is fixed.  For each
$v\in C$, the event that $v$ is an internal non-backbone vertex is the event
$v\in A\setminus B$.  Conditional on $v\in A\cup B$, this probability is
\[
  \Prb[v\in A\setminus B\mid v\in A\cup B]
  =\frac{p_v(1-p_v)}{1-(1-p_v)^2}
  =\frac{1-p_v}{2-p_v}
  \le \frac12.
\]
Moreover, these labels are independent over vertices, because the original
samples are independent over vertices and the conditioning on $A\cup B=C$ is
coordinate-wise.

Thus, along the fixed cyclic order of $T_C$, the internal non-backbone labels
are independent Bernoulli labels with probabilities at most $1/2$.  If
$M(A,B)\ge t$, then some run of $t$ consecutive positions on the tour is labeled
internal.  There are at most $n$ possible starting positions, so by a union
bound,
\[
  \Prb[M(A,B)\ge t\mid C]\le n2^{-t}.
\]
Removing the conditioning gives the same bound.

The moment estimate follows from the tail bound.  Indeed, a nonnegative random
variable satisfying $\Prb[M\ge t]\le n2^{-t}$ has $q$-th moment bounded by
$\bigl(c_0(\log n+q)\bigr)^q$ for a universal constant $c_0$.
\end{proof}

\begin{lemma}[Backbone degree bound]
There is an absolute constant $C_D$ such that
\[
  \Prb_B[Z(B)>C_D\log n]\le n^{-10}.
\]
\end{lemma}

\begin{proof}
Let $q:=\lceil 20\log_2 n\rceil$.  By Jensen's inequality,
\[
  Z(B)^q
  =\bigl(\E_A[M(A,B)\mid B]\bigr)^q
  \le \E_A[M(A,B)^q\mid B].
\]
Taking expectation over $B$ and using the previous lemma,
\[
  \E_B[Z(B)^q]
  \le \E[M(A,B)^q]
  \le \bigl(c_0(\log n+q)\bigr)^q
  \le \bigl(C'\log n\bigr)^q
\]
for an absolute constant $C'$.  Markov's inequality gives
\[
  \Prb_B[Z(B)>C_D\log n]
  \le \left(\frac{C'}{C_D}\right)^q.
\]
Choosing $C_D$ sufficiently large makes the right-hand side at most
$n^{-10}$.
\end{proof}

Therefore, with high probability over $B$, the averaged solution $x^B$ satisfies
\[
  x^B(\delta^+(b))=x^B(\delta^-(b))\le C_D\log n
  \qquad \forall b\in B\cup\{r\}.
\]

\section{Cost bound}

\begin{lemma}[Cost for a fixed pair]
For every fixed pair $(A,B)$,
\[
  \cost(y^{A,B})\le M(A,B)\,\cost(T_{A\cup B}).
\]
\end{lemma}

\begin{proof}
Consider one segment
\[
  P_j=(b_j,u_{j,1},\ldots,u_{j,k_j},b_{j+1}),
\]
and let $\ell_j$ be its length as a subpath of $T_{A\cup B}$.  By the directed
triangle inequality, for every internal vertex $u_{j,t}$,
\[
  d(b_j,u_{j,t})+d(u_{j,t},b_{j+1})\le \ell_j.
\]
Therefore the star arcs through the internal vertices of this segment have total
cost at most $k_j\ell_j$.  Also,
\[
  d(b_j,b_{j+1})\le \ell_j,
\]
so the repair arc $b_j\to b_{j+1}$ has cost at most
$(M(A,B)-k_j)\ell_j$.  Thus the total contribution of this segment is at most
\[
  k_j\ell_j+(M(A,B)-k_j)\ell_j=M(A,B)\ell_j.
\]
Summing over all segments gives
\[
  \cost(y^{A,B})\le M(A,B)\sum_j\ell_j=M(A,B)\cost(T_{A\cup B}).
\]
\end{proof}

\begin{lemma}[Expected cost]
There is an absolute constant $C_c$ such that
\[
  \E_B[\cost(x^B)]\le C_c\log n\cdot \OPT_{\mathrm{post}}.
\]
\end{lemma}

\begin{proof}
By the previous lemma,
\[
  \E_B[\cost(x^B)]
  =\E_{A,B}[\cost(y^{A,B})]
  \le \E_{A,B}\bigl[M(A,B)\cost(T_{A\cup B})\bigr].
\]
Condition on $C=A\cup B$.  The tour $T_C$ and its cost are fixed under this
conditioning.  By the longest run bound,
\[
  \E[M(A,B)\mid C]\le C''\log n.
\]
Hence
\[
  \E_{A,B}\bigl[M(A,B)\cost(T_{A\cup B})\bigr]
  \le C''\log n\cdot \E_{A,B}[\cost(T_{A\cup B})].
\]
Finally,
\[
  \cost(T_{A\cup B})\le \OPT(A)+\OPT(B),
\]
because one may concatenate optimal tours for $A$ and $B$ at the depot $r$ and
then shortcut.  Therefore
\[
  \E_{A,B}[\cost(T_{A\cup B})]
  \le 2\OPT_{\mathrm{post}}.
\]
Combining the inequalities proves the claim.
\end{proof}

\section{Conclusion for a boosted backbone}

We now finish the argument for the actual backbone used by the algorithm.  The
point is that the larger backbone must contain at least one ordinary active
sample for which the two estimates above hold.

First record the consequence of the preceding two sections.

\begin{lemma}[A good ordinary sample]
\label{lem:good-ordinary-sample}
There are absolute constants $D_0,K_0>0$ such that, if $B$ is one ordinary
active sample, then with probability at least $2/3$ over $B$, the vector $x^B$
is feasible for $\mathrm{LP}_{D_0\log n}(B)$ and satisfies
\[
  \cost(x^B)\le K_0\log n\cdot \OPT_{\mathrm{post}}.
\]
\end{lemma}

\begin{proof}
By the backbone degree bound, after increasing the absolute constant if
necessary,
\[
  \Prb_B\bigl[ Z(B)>D_0\log n\bigr]\le \frac1{12}.
\]
By the expected cost lemma and Markov's inequality, after increasing $K_0$ if
necessary,
\[
  \Prb_B\bigl[\cost(x^B)>K_0\log n\cdot \OPT_{\mathrm{post}}\bigr]
  \le \frac1{12}.
\]
Therefore both events fail with probability at most $1/6$, so both desired
properties hold with probability at least $5/6$.  In particular they hold with
probability at least $2/3$.
\end{proof}

Next we show that any larger backbone inherits feasibility from such a good
ordinary sample.

\begin{lemma}[Monotonicity under enlarging the backbone]
\label{lem:monotonicity}
Let $B\subseteq S\subseteq V$.  If $x^B$ is feasible for
$\mathrm{LP}_{D}(B)$ and $D\ge 1$, then the same vector $x^B$ is feasible for
$\mathrm{LP}_{D}(S)$.  Its cost is unchanged.
\end{lemma}

\begin{proof}
Let $U_B:=V\setminus B$ and $U_S:=V\setminus S$.  Since $B\subseteq S$, we have
$U_S\subseteq U_B$.  The support constraint for $\mathrm{LP}_{D}(B)$ says that
$x^B_{uv}=0$ for all $u,v\in U_B$; hence certainly $x^B_{uv}=0$ for all
$u,v\in U_S$.

For every vertex $v\in U_S$, we also have $v\notin B$, so the construction of
$x^B$ gives
\[
  x^B(\delta^+(v))=x^B(\delta^-(v))=p_v.
\]
Thus all outside-vertex lower degree constraints for $\mathrm{LP}_{D}(S)$ are
satisfied.

Now consider a backbone vertex $s\in S\cup\{r\}$.  If $s\in B\cup\{r\}$, then
its degree is at most $D$ by feasibility for $\mathrm{LP}_{D}(B)$.  If
$s\in S\setminus B$, then $s$ was an outside vertex with respect to the
backbone $B$, and hence
\[
  x^B(\delta^+(s))=x^B(\delta^-(s))=p_s\le 1\le D.
\]
The global balance constraints are unchanged.  Therefore all constraints of
$\mathrm{LP}_{D}(S)$ hold, and the objective value is the same vector cost.
\end{proof}

We now specify the boosted sampling.  Let $m$ be a positive integer.  Take
independent ordinary active samples
\[
  B^1,B^2,\ldots,B^m,
\]
and define the boosted backbone by
\[
  S:=\bigcup_{i=1}^m B^i.
\]
Equivalently, the vertices of $S$ are included independently with probabilities
\[
  q_v=1-(1-p_v)^m.
\]

\begin{theorem}[Degree-constrained LP for the boosted backbone]
\label{thm:boosted-backbone-lp}
There are absolute constants $a,D_0,K_0>0$ such that the following holds.  Let
$m\ge a\log n$, and let
\[
  S=\bigcup_{i=1}^m B^i
\]
be the union of $m$ independent ordinary active samples.  Then, with probability
at least $1-n^{-10}$ over the boosted backbone $S$, the LP
$\mathrm{LP}_{D_0\log n}(S)$ has a feasible solution $x$ satisfying
\[
  \cost(x)\le K_0\log n\cdot \OPT_{\mathrm{post}}.
\]
\end{theorem}

\begin{proof}
For each $i$, apply Lemma~\ref{lem:good-ordinary-sample} to $B^i$.  The events
\[
  G_i:=\{x^{B^i}\text{ is feasible for }\mathrm{LP}_{D_0\log n}(B^i)
  \text{ and }\cost(x^{B^i})\le K_0\log n\cdot \OPT_{\mathrm{post}}\}
\]
are independent and each has probability at least $2/3$.  Hence
\[
  \Prb[\text{none of }G_1,\ldots,G_m\text{ occurs}]
  \le \left(\frac13\right)^m.
\]
Choosing the absolute constant $a$ sufficiently large gives
\[
  \left(\frac13\right)^m\le n^{-10}.
\]
On the complementary event, choose an index $i$ for which $G_i$ occurs.  Since
$B^i\subseteq S$, Lemma~\ref{lem:monotonicity} implies that the same vector
$x^{B^i}$ is feasible for $\mathrm{LP}_{D_0\log n}(S)$ and has the same cost.
This proves the theorem.
\end{proof}

\begin{remark}[Sampling with probabilities $\min\{1,mp_v\}$]
The same conclusion also holds if the boosted backbone is sampled independently
with probabilities
\[
  q'_v:=\min\{1,mp_v\}.
\]
Indeed,
\[
  1-(1-p_v)^m\le \min\{1,mp_v\}=q'_v,
\]
so this distribution stochastically dominates the union-of-$m$-samples
distribution.  Using one uniform random variable per vertex, one may couple a
set $S_0$ distributed as $\bigcup_{i=1}^m B^i$ and a set $S'$ distributed with
marginals $q'_v$ so that $S_0\subseteq S'$ always.  Whenever
$\mathrm{LP}_{D_0\log n}(S_0)$ has the good witness from
Theorem~\ref{thm:boosted-backbone-lp}, Lemma~\ref{lem:monotonicity} transfers
that same witness to $S'$.  Thus the theorem remains valid for the more common
boosting rule $q'_v=\min\{1,mp_v\}$.
\end{remark}

\end{document}
